\documentclass[11pt,a4paper]{article}

\usepackage[utf8]{inputenc}
\usepackage[T1]{fontenc}
\usepackage[brazil]{babel}
\usepackage[margin=2.5cm]{geometry}
\usepackage{amsmath,amssymb}
\usepackage{listings}
\usepackage{xcolor}
\usepackage{booktabs}
\usepackage{tabularx}
\usepackage{xltabular}
\usepackage{tcolorbox}
\tcbuselibrary{breakable,skins}
\usepackage{enumitem}
\usepackage{hyperref}
\usepackage{lmodern}
\usepackage{parskip}

\hypersetup{colorlinks=true, linkcolor=blue!60!black, urlcolor=blue!60!black}

\definecolor{codebg}{gray}{0.96}
\definecolor{codecomment}{rgb}{0.25,0.45,0.25}
\definecolor{codekw}{rgb}{0.13,0.13,0.55}
\definecolor{codestr}{rgb}{0.55,0.25,0.25}

\lstset{
  language=Python,
  backgroundcolor=\color{codebg},
  basicstyle=\ttfamily\small,
  commentstyle=\color{codecomment}\itshape,
  keywordstyle=\color{codekw}\bfseries,
  stringstyle=\color{codestr},
  showstringspaces=false,
  breaklines=true,
  frame=single,
  framerule=0pt,
  xleftmargin=1em,
  aboveskip=1em,
  belowskip=1em,
}
\lstdefinestyle{shell}{language=bash,keywordstyle=\color{codekw}}

% ---- notacao do exercicio -------------------------------------------------
\newcommand{\Om}{\Omega_m}
\newcommand{\OL}{\Omega_\Lambda}
\newcommand{\Ok}{\Omega_k}
\newcommand{\M}{\mathcal{M}}

% ---- caixa das notas que eu ainda tenho que escrever ----------------------
\newtcolorbox{escrever}[1][Para escrever]{
  colback=yellow!6, colframe=orange!55!black, coltitle=white,
  fonttitle=\bfseries\small, title={#1},
  boxrule=0.6pt, left=6pt, right=6pt, top=4pt, bottom=4pt,
  breakable
}

% ---- marcador de lacuna --------------------------------------------------
\newcommand{\pendente}[1]{\textcolor{red!70!black}{\textbf{[#1]}}}

\title{\textbf{Notas para o exercício 1 de cosmologia}\\[0.4em]
\large Exercício 1 de Cosmologia Numérica}
\author{Fabio Fioravante \\ \small Cosmologia Numérica --- Exercício 1}
\date{Setembro de 2026}

\begin{document}
\maketitle

Estas são as notas pessoais que eu tomei enquanto aprendia as ferramentas de cosmologia e estatística que usei no Exercício 1 da disciplina
de Cosmologia Numérica, junto com os detalhes das decisões e dos resultados que ficaram de fora do README. O README,
\texttt{exercicio1-readme.tex}, é curto e traz só o necessário para usar o código e as hipóteses assumidas.

\tableofcontents

% ==========================================================================
\section{Introdução}
\label{sec:historia}

Supernovas são velas-padrão, então têm mais ou menos o mesmo brilho intrínseco por que acontecem 
mais ou menos na mesma época e condições de uma estrela. Com isso podemos saber qual a distância que elas estão, 
olhando justamente para a discrepância entre o brilho previsto e o observado para cada uma delas, e o redshift, comparando
a frequência da luz que chega e a prevista.

A partir desses dados podemos estimar os parâmetros cosmológicos que queremos, pois eles controlam o ritmo de expansão do universo, e
portanto a relação entre distância e redshift.

Para a parte de estatística que vai servir de ferramenta para interpretar os dados, vimos em aula, o teorema de bayes é 
\begin{equation}
    P(\theta|d)=\frac{P(d|\theta) P(\theta)}{P(d)},
\end{equation}
que com os nomes é algo do tipo
\begin{equation}
    \text{posteriori}=\frac{\text{likelihood}\times\text{priori}}{\text{evidência}}
\end{equation}

O que queremos é encontrar a distribuição de $\Omega_m$ e $\Omega_\Lambda$ que melhores se ajustam aos dados. A evidência é uma normalização que não vai importar aqui, por ser a mesma para todas as posteriores, então vamos descartá-la.
Basicamente, damos a priori, que é o que achamos razoável antes de olhar os dados, e a likelihood, que é a nota que cada curva recebe de 
acordo com os dados e os erros. Ademais, os erros das supernovas não são independentes, portanto, precisamos lidar com a matriz de covariância, que registra quais erros são correlacionados.

Para calcularmos a posteriori, uma forma é fazer uma grade de valores de $(\Omega_m, \Omega_\Lambda)$, e calcular a posteriori para cada ponto da grade. 
Mas isso explode quando temos mais parâmetros, então vamos usar o MCMC (Markov Chain Monte Carlo), 
que é faz um passeio aleatório pelo plano $(\Omega_m, \Omega_\Lambda)$ onde, dado um ponto inicial, 
sorteia um próximopasso, aceita sempre se a posterior melhora, e aceita só às vezes se ela piora. 
Essa regra faz o passeio passar mais tempo onde a probabilidade é maior, na proporção certa, 
e, fazendo o historgrama de pontos visitados, por vezes que cada ponto foi visitado, obtemos a posteriori.

Por fim, ainda temos que lidar com o fato de que conhecemos os brilho das supernovas a menos de uma constante, isso faz com que todas as
distâncias previstas tenham que ser ajustadas por uma constante $\M$, que o exercício pede para lidarmos.

% ==========================================================================
\section{O que o código faz}

Essa seção explica passo a passo o que cada módulo faz, como foi pensado, o que foi admitido e o que foi testado.

% --------------------------------------------------------------------------
\subsection{Por que o pacote é dividido assim}

Um \emph{módulo} em Python é um arquivo \texttt{.py} com funções e classes que
outros arquivos importam. O pacote é dividido em módulos que não sabem nada uns
dos outros, e o motivo é poder testar cada um sozinho: o amostrador é testado
contra uma distribuição cuja resposta é conhecida de antemão, sem cosmologia
envolvida, e a cosmologia é testada contra casos em que a conta sai à mão, sem
dados envolvidos. Quando o resultado final sai errado, isso diz em qual das
caixas procurar o erro.

O enunciado pede três módulos, e a leitura dos dados pode ficar dentro de
\texttt{bayes} ou num quarto módulo próprio. Ela ficou num quarto módulo,
\texttt{union21}, pelo mesmo motivo. No notebook, a Parte 3 (célula 3)
é a única que abre os arquivos do catálogo, e a priori (célula 4) também depende
dele, porque testa $E^2 > 0$ até o maior redshift da tabela. Colocar essas duas
peças em \texttt{bayes} faria o amostrador depender dos arquivos de dados, e o
teste da gaussiana, que não tem nada a ver com supernovas, passaria a precisar
deles. O nome diz de onde vêm os dados: outra compilação de supernovas seria
outro módulo com as mesmas duas funções, \texttt{log\_likelihood(theta, model)}
e \texttt{log\_prior(theta)}, e \texttt{bayes} não precisaria mudar.

\begin{escrever}[Para escrever --- nota própria]
Dar um exemplo de erro no código que o teste da gaussiana (seção~\ref{sec:bayes})
pegaria e o teste da cosmologia não pegaria, e um exemplo do contrário.
\end{escrever}

% --------------------------------------------------------------------------
\subsection{\texttt{cosmology}}
\label{sec:cosmo}

Definimos, no primeiro módulo do código, as funções cosmológicas que usaremos na classe \texttt{FLRW}.

Teoricamente, dada a métrica FLRW e as equações de Friedmann, o ritmo de expansão do universo é dado por
\begin{equation}
E^2(z)= \frac{H^2(z)}{H_0^2} = \Om(1+z)^3 + \Omega_r(1+z)^4 + \Ok(1+z)^2 + \OL,
\end{equation}
onde $\Omega_r$ é a densidade de radiação, que vamos desprezar, e $\Ok = 1 - \Om - \OL$ é a densidade de curvatura. A partir do ritmo de expansão, 
podemos calcular a distância comóvel $\chi(z)$, que é o caminho percorrido pela luz desde a emissão até hoje, em unidades de $c/H_0$:


$$
\frac{1}{E(z)}=\frac{d\chi(z)}{dz}\to \int_0^z \frac{dz'}{E(z')}=\chi(z).
$$

A distância comóvel radial é então dada por

$$
r(z)=\frac{c}{H_0}\chi(z)=\frac{c}{H_0}\int_0^z \frac{dz'}{E(z')},
$$

e o $S_k$ (não sei como posso chamar) é uma função da distância comóvel radial e pode ser escrito, para uma curvatura $k$ arbitrária, como:

\begin{equation}
S_k(r)=\begin{cases}
R_0 \sin (r/R_0)\ \ \ \ (k=+1)    \\
r \ \ \ \ (k=0)    \\
R_0 \sinh (r/R_0)\ \ \ \ (k=-1),
\end{cases}
\end{equation}

com $R_0$ sendo o raio de curvatura do universo, e pode ser escrito como função de $\Omega_k$ quando olhamos para a parte de curvatura da equação de Friedmann:

\begin{equation}
1-\Omega_{tot}=-\frac{k c^2}{a^2 R_0^2 H^2}
\end{equation}

Avaliando a equação de curvatura hoje ($a_0 = 1$, $H = H_0$) e definindo

$\Omega_k \equiv 1 - \Omega_{\rm tot,0}$:

\begin{equation}
\Omega_k = -\frac{k c^2}{R_0^2\, H_0^2}
\quad\Longrightarrow\quad
R_0 = \frac{c/H_0}{\sqrt{|\Omega_k|}} = \frac{D_H}{\sqrt{|\Omega_k|}},
\end{equation}

com $D_H \equiv c/H_0$. O sinal de $\Omega_k$ é oposto ao de $k$:

$k=+1$ (fechado) $\Leftrightarrow \Omega_k < 0$, e $k=-1$ (aberto) $\Leftrightarrow \Omega_k > 0$.

Como $r(z)=D_H\chi(z)$, o argumento que entra em $S_k$ pode ser escrito como

$$
\frac{r}{R_0}=\frac{D_H\chi}{R_0}
=\sqrt{|\Omega_k|}\,\chi(z),
$$

e portanto:

\begin{equation}
S_k(r(z)) =
\begin{cases}
\dfrac{c}{H_0\sqrt{ \Omega_k}}\,\sinh\!\left(\sqrt{\Omega_k}\;\chi(z)\right), & \Omega_k > 0\\[10pt]
\frac{c}{H_0}\,\chi(z), & \Omega_k = 0\\[10pt]
\dfrac{c}{H_0 \sqrt{|\Omega_k|}}\,\sin\!\left(\sqrt{|\Omega_k|}\;\chi(z)\right), & \Omega_k < 0.
\end{cases}
\end{equation}

e, por fim, podemos definir as distâncias de luminosidade e angular, respectivamente:
\begin{equation}
d_L=S_k (r=\chi\frac{c}{H_0})(1+z)
\end{equation}
e
\begin{equation}
d_A=S_k (r=\chi\frac{c}{H_0})/(1+z)
\end{equation}

Para escalas astronômicas, é comum usar o módulo de distância, que é uma escala logarítmica da distância de luminosidade:
\begin{equation}
  \mu= 5\log_{10}\left(\frac{d_L}{\rm Mpc}\right)+25,
\end{equation}
que é também o dado que temos na tabela do Union2.1.

Além das distâncias, a classe calcula o tempo $t(z)$ e o elemento de de volume $dv/dz\,d\Omega$, que o enunciado também pede.

Por sinal, usei o livro da barbara ryden como referencia nessa parte toda

\paragraph{testes}

Os testes pedidos para essa primeira parte são 
\begin{itemize}
  \item E(0)=1, checar se hoje $H=H_0$;
  \item Lei de Hubble para z baixo, $D_L \to cz/H_0$ quando $z \to 0$;
  \item $D_M = (c/H_0)\,\chi$ quando $\Ok = 0$;
  \item Checar a relação entre as distâncias, $D_L = (1+z)^2 D_A$ para qualquer curvatura.
\end{itemize}
Todos esses testes são feitos na primeira célula (\texttt{\#testes1}) e todos passam 


% --------------------------------------------------------------------------
\subsection{\texttt{estatística}}
\label{sec:bayes}

Retomando um pouco do que foi discutido na introdução, como ferramenta para analisarmos os dados, vamos usar a abordagem bayesiana para probabilidade.

Nessa abordagem, muda-se a interpretação de probabilidade que, usualmente, é entendida como o número de eventos dividido pelo número de tentativas.
Essa interpretação é chamada de frequentista, e é a mais comum, contudo, ela depende de fenômenos que tenham caráter repetitivo, o que não
é o caso para todos os eventos, existem certas perguntas que a abordagem frequentista não consegue responder, como por exemplo, 
qual a probabilidade de ter chovido no dia 14 de setembro de 2026. A abordagem Bayesiana interpreta a probabilidade como o grau de confinça, logicamente deduzido
dada a informação que se tem. Essa abordagem é a que vamos usar para estimar os parâmetros cosmológicos, dado o conjunto de dados do Union2.1.

Dessa noção de probabilidade um dos primeiros resultados que podemos tirar, e o que vai estruturar tudo que faremos daqui para frente, é o teorema
de Bayes, onde, no nosso contexto, ele pode ser escrito como
\begin{equation}
  p(\theta|d)=\frac{p(d|\theta)p(\theta)}{p(d)},
\end{equation}
onde $\theta$ é o conjunto de parâmetros, aqui $(\Om,\OL)$, e $d$ são os dados observados.

Cada fator dessa relação tem um nome no linguajar bayesiano e eles são:
\begin{itemize}
  \item \textbf{posterior} $p(\theta|d)$: é a probabilidade de $\theta$ dado os dados, ou seja, é o que queremos estimar;
  \item \textbf{sampling} $p(d|\theta)$: é a probabilidade de observar os dados dado o valor de $\theta$;
  \item \textbf{priori} $p(\theta)$: é a probabilidade de $\theta$ antes de observar os dados;
  \item \textbf{evidência} $p(d)$: é a probabilidade dos dados, independentemente do valor de $\theta$,
\end{itemize}
e considerando o sampling com dados fixos e função dos parâmetros temos a função de verossimilhança, ou likelihood
\begin{equation}
  \mathcal{L}(\theta) = p(d|\theta).
\end{equation}

Colocando o linguajar que definimos, o teorema de bayes é 
\begin{equation}
    \text{posteriori}=\frac{\text{likelihood}\times\text{priori}}{\text{evidência}}
\end{equation}

Para calcular a posteriori para os parâmtros que queremos ($\Om$, $\OL$), podemos fazer uma amostragem via MCMC. 
Isso que dizer que, dado um ponto inicial, um par $(\Omega_{m_0}, \Omega_{\Lambda_0})$, sorteamos o próximo ponto de manera aletória, dependendo apenas do ponto atual,
e aceitamos um ponto que melhora a posteriori, e aceitamos um ponto que piora a posteriori com uma probabilidade proporcional 
à razão entre as duas posteriores. Formalmente, isso é o que o algorimto Metropolis-Hastings faz, e é o que vamos usar para amostrar a posteriori.
Um ponto importante desse algoritmo que vale citar é que caso estejamos em um ponto e o próximo ponto sorteado seja recusado, o 
ponto atual é anotado de novo, e isso faz com que as regiões com maior posteriori sejam visitadas mais vezes, além disso, também
é importante que possam ser escolhidos pontos que pioram a posteriori, pois isso faz com que a cadeia explore a distribuição inteira,
e não fique presa em um pico local de posteriori.

No código, definimos 4 classes: \texttt{Posterior}, \texttt{Proposal}, \texttt{GaussianRandomWalk} e \texttt{MetropolisHastings}.
De maneira geral, \texttt{Posterior} recebe a priori e a likelihood e, para um $\theta$, devolve o logaritmo da posterior. 
Se $\theta$ estiver fora do intervalo permitido pela priori, ela devolve $-\infty$, que é o logaritmo de zero, ou seja, probabilidade zero. 
A \texttt{Proposal} é uma lista de regras que o programa precisa seguir para sortear novos pontos, que são 
sortear (\texttt{propose}) e dizer qual era a probabilidade de sortear um dado salto (\texttt{logpdf}).
O \texttt{GaussianRandomWalk} é uma regra para sortear novos pontos, que é o ponto atual mais um ruído gaussiano. Como
não podemos afirmar que a posteriori é simétrica, precisamos deformar o ruído no 
formato da posteriori, e isso é feito via decomposição de Cholesky da matriz de covariância desejada,
$\Sigma=L L^T$.
Por fim, a classe \texttt{MetropolisHastings} é o laço da cadeia, que recebe o gerador de números aleatórios de fora, como argumento, para que a mesma semente (o número que inicializa o gerador) produza exatamente a mesma cadeia. Ele devolve as amostras, o log da posterior em cada uma, quais passos foram aceitos e a taxa de aceitação.


Tudo é feito com logaritmos. A verossimilhança de 580
pontos é o produto de 580 números pequenos, e um produto desses facilmente fica
menor que o menor número que o computador representa,
virando zero. Com logaritmos não temos esse problema, pois todos os termos se somam. 
A regra de
aceitação fica
\begin{equation}
\log u < \log\pi(\theta') - \log\pi(\theta)
+ \underbrace{\log q(\theta\,|\,\theta') - \log q(\theta'\,|\,\theta)}_{\text{correção de Hastings}},
\end{equation}
onde $u$ é um número aleatório uniforme entre 0 e 1, $\pi$ é a posterior, $\theta$ o ponto atual, $\theta'$ o candidato e
$q(\theta'|\theta)$ a probabilidade de sortear $\theta'$ estando em $\theta$. A
correção de Hastings compensa regras de sorteio que preferem uma direção.
Para o passeio gaussiano, ir de $\theta$ a $\theta'$ é tão provável quanto
voltar, os dois termos se cancelam e a correção vale zero. O código calcula a
correção mesmo assim, via \texttt{logpdf}, para poder trocar de regra de sorteio
depois sem mexer no laço.

O tamanho do passo decide a eficiência. Passos pequenos demais são quase sempre
aceitos, mas a cadeia anda devagar e cada amostra repete quase a anterior.
Passos grandes demais caem fora da região provável e são quase sempre
recusados. Uma receita clássica para começar, válida para alvos parecidos com
gaussianas, é usar como covariância da proposta a covariância da posterior
multiplicada por $2{,}38^2/d$, onde $d$ é o número de parâmetros. A taxa de
aceitação é só um sintoma, e o que realmente mede a eficiência são os
diagnósticos da seção~\ref{sec:diag}.

\paragraph{testes} Os testes pedidos para essa segunda parte são
\begin{itemize}
  \item Usar o algoritmo para recuperar uma gaussiana, dados os desvios
  \item Teste de reproduzibilidade, rodar duas vezes a mesma cadeia com a mesma semente e ver se os resultados batem
  \item Testar se um estado rejeitado é repetido
\end{itemize}

Todos esses testes são feitos na segunda célula (\texttt{\#testes2}) e todos passam.

% --------------------------------------------------------------------------
\subsection{\texttt{análise do Union2.1}}
\label{sec:like}

Definimos, para cada supernova, o resíduo, como sendo, a diferença entre o $\mu$ (módulo de distância) observado e 
o $\mu$ previsto. Se os erros fossem independentes, o $\chi^2$ seria apenas a soma de
cada $\chi^2$ para cada supernova. Como a matriz de covariância não é diagonal, os erros das supernovas
não são independentes, precisamos tratar isso de maneira diferente.

Nesse caso o $\chi^2$ é dado por
\begin{equation}
  \chi^2=(\mu_{\rm obs}-\mu_{\rm th})^T C^{-1}(\mu_{\rm obs}-\mu_{\rm th}),
\end{equation}
onde $C$ é a matriz de covariância. E a maneira de fazemos isso que exige menos numericamente (não sei o por quê exatamente)
é fatorar a matriz de covariância via Cholesky, $C=LL^T$, onde $L$ é uma matriz
triangular inferior e resolver o sistema $Ly = (\mu_{\rm obs}-\mu_{\rm th})$ e $\chi^2=y^Ty$.

Fazendo isso no código, começamos na célula 3 do notebook, lendo os dados 
e checando tamanho, simetria, etc. Só então fatora a covariância:
\begin{lstlisting}
L_sys = np.linalg.cholesky(cov_sys)           

uns = np.ones(N_SN)
w = solve_triangular(L_sys, uns, lower=True)   # resolve L_sys @ w = 1
C1 = w @ w
\end{lstlisting}
Como $C$ não depende de $\Om$ nem de
$\OL$, basta fazê-la uma vez para todos os pontos da cadeia. As duas últimas linhas preparam dois ingredientes
da marginalização de $\M$ que também não dependem dos parâmetros, $w$, a solução
de $Lw = \mathbf{1}$ com $\mathbf{1}$ o vetor de 580 uns, e $C_1 = w^Tw$
(seção~\ref{sec:M}).

A cada ponto que a cadeia visita, entram duas funções:
\begin{lstlisting}
def chi2_marg(mu_teoria):
    delta = mu_obs - mu_teoria
    y = solve_triangular(L_sys, delta, lower=True)   # resolve L_sys @ y = delta
    A = y @ y
    B = y @ w
    return A - B**2 / C1

def log_likelihood(theta, model):
    omega_m, omega_lambda = theta
    with np.errstate(invalid="ignore", divide="ignore"):
        cosmo = model(H0=70, omega_m=omega_m, omega_lambda=omega_lambda)
        mu_teoria = cosmo.mu(z_sn)
    if not np.all(np.isfinite(mu_teoria)):
        return -np.inf
    return -0.5 * chi2_marg(mu_teoria)
\end{lstlisting}

A primeira, \texttt{chi2\_marg}, recebe os $\mu$ teóricos das 580 supernovas,
forma o resíduo \texttt{delta} e resolve $Ly = \Delta$ com a função
\texttt{solve\_triangular} do scipy. O argumento \texttt{lower=True} avisa que
$L$ é triangular inferior. Resolver $Ly = \Delta$ dá o
mesmo $\chi^2$ da equação acima porque
$C^{-1} = (LL^T)^{-1} = (L^{-1})^T L^{-1}$, e então
$\Delta^T C^{-1} \Delta = (L^{-1}\Delta)^T (L^{-1}\Delta) = y^Ty$, sem que
$C^{-1}$ precise ser calculada em momento nenhum. Esse $y^Ty$ é o \texttt{A} da
função, e o que ela devolve, $A - B^2/C_1$, é o mesmo $\chi^2$ já com $\M$
integrado, como explicado na seção~\ref{sec:M}.

A segunda, \texttt{log\_likelihood}, tem a assinatura que a classe
\texttt{Posterior} espera, recebe o ponto \texttt{theta} e a classe cosmológica
\texttt{model}, que na prática é a \texttt{FLRW}. Ela seta um universo com
$H_0 = 70$, e
calcula $\mu$ para as 580 supernovas. O \texttt{np.errstate} só silencia os avisos que o numpy dá quando alguma
conta vira \texttt{nan}. Se algum $\mu$ não é um número finito, que é o caso da
luz que passou do ponto oposto de um universo fechado, a função devolve
$-\infty$, o logaritmo de zero, e o ponto é recusado. No caso normal ela devolve
$-\chi^2/2$, porque a verossimilhança gaussiana é proporcional a
$e^{-\chi^2/2}$ e a cadeia trabalha com logaritmos.



Um último ponto para se consederar é que, para algumas combinações de $(\Om,\OL)$, com muita energia escura e pouca
matéria, $E^2(z)$ fica negativo antes de chegar ao maior redshift do catálogo
($z = 1{,}414$). Nesses universos, voltando no tempo, a expansão para e se
inverte antes de chegar a esse $z$, e não houve Big Bang. Uma supernova com
aquele redshift simplesmente não poderia existir, então o ponto é impossível,
e o código tem que recusá-lo de propósito na priori. Deixar a raiz quadrada de
um número negativo virar \texttt{nan} e ser recusada por acidente funciona, mas
esconde a decisão. Um exemplo: $(\Om,\OL) = (0{,}1;\,2{,}0)$ tem $E^2 < 0$ dentro
da faixa do catálogo.







\end{document}
